Read in Gauge Theory Gravity: the metric $g$ below is the derived metric of the gravitational field $(\hb, \Omega)$ on flat spacetime, the manifold is the invertible region, and the tractor connection is the flat-space $\Cl(2,4)$ gauge field of Theorem \claimref{synthesis.gtg-tractor-gauge-field-reproduces-tractor-connection}; the index computation below is its reading in the splitting of $g$ (Fact \claimref{synthesis.imported-gtg-dictionary}).

\emph{Setting.} The tractor form in the splitting of $g$ is $\ip{(\sigma,\mu,\rho)}{(\sigma',\mu',\rho')} = \sigma\rho' + \rho\sigma' + g^{bc}\mu_b\mu'_c$, preserved by the connection (Lemma \claimref{interior.metric-parallelism-trace-schouten-condition}(i)). By Fact \claimref{interior.imported-tractor-connection}, $\nabla_v(\sigma,\mu,\rho) = (\nabla_v\sigma - \mu(v),\ \nabla_v\mu + v^\flat\rho - \Sch(v,\cdot)\sigma,\ \nabla_v\rho + \Sch(v,\mu^\sharp))$. For the constant sections $I_g = (1,0,0)$ and $X = (0,0,1)$ this gives $\nabla_v I_g = (0, -\Sch(v,\cdot), 0)$, which is Theorem \claimref{interior.covariant-form-vacuum-equations}, and $\nabla_v X = (0, v^\flat, 0)$.

\emph{(1) Shape.} Since $\ip{(\sigma,\mu,\rho)}{I_g} = \rho$, the orthogonal complement of $I_g$ is $\{\rho = 0\}$, and a totally isotropic $W \ni I_g$ lies in it. Subtracting multiples of $I_g$ from two further basis vectors removes their $\sigma$ components, so $W = \mathrm{span}(I_g, \ell, m)$ with $\ell, m$ in the $\mu$ slot, and isotropy of $W$ is exactly total nullity of $A_W = \mathrm{span}(\ell, m) \subset T_xM\otimes\C$. The same argument with $\ip{(\sigma,\mu,\rho)}{X} = \sigma$ handles point vacua. A totally null plane $A$ in four dimensions equals its own orthogonal complement. If $A \cap \bar A = 0$, then $T_xM\otimes\C = A \oplus \bar A$ and every real vector is $a + \bar a$ with $a \in A$, with $g(a + \bar a, a + \bar a) = 2g(a, \bar a)$; the Hermitian form $h(a,b) = g(a,\bar b)$ on the two-dimensional $A$ would then give $g$ a signature $(2p, 2q)$ on the real space, contradicting $(1,3)$. So $A \cap \bar A \neq 0$, and it is conjugation-invariant, so it contains a nonzero real vector $\ell_W$, null since $A$ is totally null. Two real null vectors in $A$ are orthogonal, and orthogonal real null vectors in signature $(1,3)$ are proportional, so the real null direction is unique. Conversely, for a real null $\ell$ the quotient $\ell^\perp/\ell$ is two-dimensional with definite induced form, whose complexification has exactly two null lines; their preimages are the two totally null planes containing $\ell$, so $\ell$ lies in exactly two infinity vacua. The identification with Fock vacua is the purity statement of Fact \claimref{one-relation.imported-pure-spinors-incidence-chart}: each conformal spinor $u$ has annihilator $W_u = \{v : c(v)u = 0\}$ of dimension three, determining $u$ up to scale, so for $u \in \CS$, $I_g \in W_u$ if and only if $c(I_g)u = 0$; the kernel of $c(I_g)$ is four-dimensional on the eight-dimensional Fock module: $c(I_g)^2 = Q(I_g) = 0$ gives $\mathrm{im}\,c(I_g) \subseteq \ker c(I_g)$, and $c(I_g)c(X) + c(X)c(I_g) = 2\ip{I_g}{X} = 2$ gives $\psi = \tfrac12 c(I_g)c(X)\psi \in \mathrm{im}\,c(I_g)$ for $\psi \in \ker c(I_g)$, so kernel and image coincide and each has half the dimension. The two totally null planes through a real null direction belong to the two families (two maximal isotropic subspaces of $\C^6$ lie in the same family exactly when their intersection has odd dimension, and the two infinity vacua through $\ell$ meet in $\mathrm{span}(I_g, \ell)$, of dimension two).

\emph{(2) Classes.} For $w' \in W$ and a function $f$, $\ip{\nabla_v(f w)}{w'} = f\ip{\nabla_v w}{w'} + (vf)\ip{w}{w'} = f\ip{\nabla_v w}{w'}$ because $W$ is isotropic, so $\omega_v$ is tensorial in $w$, and by $\ip{\nabla_v w}{w'} = -\ip{w}{\nabla_v w'}$ in $w'$ as well; it depends only on $W$ at $x$. With $w = I_g$ and $w' = (0, \mu', 0)$: $\omega_v(I_g, w') = g^{bc}(-\Sch(v,\cdot))_b\mu'_c = -\Sch(v, \mu'^\sharp)$, and $\omega_v(I_g, I_g) = 0$. With $w = X$: $\omega_v(X, w') = g(v, \mu'^\sharp)$ and $\omega_v(X, X) = 0$. For the Fock statement (with $W$ in the family of $\CS$), $c(\nabla_v I_g)u_W = 0$ if and only if $\nabla_v I_g \in W_{u_W} = W$; a $\mu$-slot vector lies in $W$ exactly when it lies in $A_W$, and since $A_W = A_W^\perp$, exactly when it is orthogonal to $\ell$ and $m$, that is when $\omega_v(I_g, \ell) = \omega_v(I_g, m) = 0$.

\emph{(3) Point vacua.} $\omega_{\ell_W}(X, w') = g(\ell_W, w')$ vanishes on $A_W$ because $A_W$ is totally null and contains $\ell_W$; no curvature enters.

\emph{(4) Einstein.} (b) implies (a) trivially. For (a) implies (c): by (1) every real null direction is some $\ell_W$, so (a) says $\Sch(\ell,\ell) = 0$ for all real null $\ell$. The symmetric bilinear forms vanishing on the null cone form the line spanned by $g$: in an orthonormal frame, $\ell = e_0 \pm e_i$ give $\Sch_{00} + \Sch_{ii} = 0$ and $\Sch_{0i} = 0$, and $\ell = e_0 + (\cos\theta) e_i + (\sin\theta) e_j$ with $\cos\theta, \sin\theta$ both nonzero then gives $\Sch_{ij} = 0$ for $i \neq j$; so $\Sch = c\,g$ with $c = \Sch_{00} = -\Sch_{ii}$, and taking the trace $c = J/4$. Lemma \claimref{interior.metric-parallelism-trace-schouten-condition}(ii) gives $c = R/24$ and $\Ric = \tfrac{R}{4}g$. For (c) implies (b): if $\Sch = \tfrac{J}{4} g$ then $\omega_{\ell_W}(I_g, w) = -\tfrac{J}{4}g(\ell_W, w) = 0$ for $w \in A_W$ by total nullity. Constancy of $R$ on a connected region where $\Ric = \tfrac{R}{4}g$: the contracted Bianchi identity $\nabla^a\Ric_{ab} = \tfrac12\nabla_b R$ (an identity of the Levi-Civita connection, unchanged by the overall sign convention for $\Ric$ since both sides change sign together) gives $\tfrac14\nabla_b R = \tfrac12\nabla_b R$, so $\nabla R = 0$ and $R$ is constant on the connected region (Schur's argument in dimension four); then $J = R/6$, $\Sch = \tfrac{\Lambda}{6}g$ with $\Lambda = R/4$, and $\omega_v(I_g, w) = -\tfrac{\Lambda}{6}g(v,w)$ by (2).

\emph{(5) Ricci-flat.} If every infinity vacuum has vanishing $(I_g,\cdot)$ classes along every $v$, then $\Sch(v, \ell) = 0$ for every $v$ and every real null $\ell$, and real null vectors span $T_xM$, so $\Sch = 0$; conversely $\Sch = 0$ makes every class $-\Sch(v,\cdot)$ vanish. $\Sch = 0$ if and only if $\Ric = 0$ by Lemma \claimref{interior.metric-parallelism-trace-schouten-condition}(ii), and if and only if $\nabla I_g = 0$ by Theorem \claimref{interior.covariant-form-vacuum-equations}.

\emph{(6) One vacuum.} The condition is $\Sch(v, a) = 0$ for all $v$ and all $a \in A = A_W$: the map $\Sch : T_xM\otimes\C \to (T_xM\otimes\C)^*$ kills $A$, so by symmetry its image annihilates $A$, hence lies in the annihilator $A^\flat$ (using $A^\perp = A$), and $\Sch \in \mathrm{Sym}^2 A^\flat$. Since $\Sch$ is real it also lies in $\mathrm{Sym}^2 \bar A^\flat$. For two planes, $\mathrm{Sym}^2 A^\flat \cap \mathrm{Sym}^2 \bar A^\flat = \mathrm{Sym}^2 (A \cap \bar A)^\flat$ (a symmetric tensor in $\mathrm{Sym}^2U \cap \mathrm{Sym}^2U'$ has all its contractions in $U \cap U'$), and $A \cap \bar A = \C\ell_W$ by (1). Hence $\Sch = c\,\ell_W^\flat\otimes\ell_W^\flat$, $c$ real. Conversely such $\Sch$ has $\Sch(v, a) = c\,g(v,\ell_W)g(\ell_W, a) = 0$ for $a \in A$. Its trace is $c\,g(\ell_W,\ell_W) = 0$, so $J = R/6 = 0$ and $\Ric = 2\Sch + \tfrac{R}{6}g = 2c\,\ell_W^\flat\otimes\ell_W^\flat$.

\emph{(7) Matter.} For null $\ell$, $\Sch(\ell,\ell) = \tfrac12\Ric(\ell,\ell)$, so the condition $-\Sch(\ell,\ell) = 4\pi G\,T(\ell,\ell)$ reads $(\Ric + 8\pi G\,T)(\ell,\ell) = 0$ for every real null $\ell$, and the argument of (4) applied to the real symmetric form $\Ric + 8\pi G\,T$ gives $\Ric + 8\pi G\,T = h\,g$; the converse is immediate. Now let this hold on a connected region with $\nabla^aT_{ab} = 0$, and write $\tau := g^{ab}T_{ab}$. The trace gives $R + 8\pi G\,\tau = 4h$. The divergence gives $\nabla^a\Ric_{ab} = \nabla_b h$, and the contracted Bianchi identity gives $\nabla^a\Ric_{ab} = \tfrac12\nabla_b R = \tfrac12\nabla_b(4h - 8\pi G\tau) = 2\nabla_b h - 4\pi G\nabla_b\tau$. Hence $\nabla_b h = 4\pi G\,\nabla_b\tau$, and $\Lambda := h - 4\pi G\tau$ is constant. For the sign remark: let $\Ric^{\circ}$ be the Ricci tensor of the convention in which the round sphere has positive Ricci curvature; it is unchanged under $g \mapsto -g$, and de Sitter space has $\Ric^{\circ} = -\Lambda g$ in signature $(1,3)$, so the store's normalisation $\Ric = \Lambda g$ means $\Ric = -\Ric^{\circ}$. The Einstein equations $\Ric^{\circ} = 8\pi G(T - \tfrac12 \tau g) - \Lambda g$, written in signature $(1,3)$, then read $\Ric + 8\pi G\,T = (\Lambda + 4\pi G\,\tau)\,g$, which is the form above.

\emph{The check.} The script verifies, with exact Gaussian-rational arithmetic, that the non-Levi part of the connection preserves the tractor form, the class formulas of (2) and (3) on a generic rational $\Sch$, the own-ray system of (4) on eleven rational null vectors (rank nine on the ten-dimensional space of symmetric forms, kernel spanned by $g$, unchanged when the $(I_g, m)$ equations are added) and the transverse classes $-\tfrac{\Lambda}{6}g(v,w)$ at $\Lambda = 3$, the one-vacuum kernel of (6) (spanned by $\ell^\flat\otimes\ell^\flat$, trace-free), the Fock criterion of (2) in the exterior model of $\Cl(2,4)_\C$ (the vacuum of $W$ is annihilated by $c(\nabla_v I_g)$ exactly when both classes vanish, for a generic, a radiation and an Einstein $\Sch$), the four-dimensional kernels of $c(I_g)$ and $c(X)$ with zero intersection on the eight-dimensional module, the real-direction count of (1) on sample planes, and for (7), with $\kappa := 4\pi G$ and $\Lambda$ rational and a rational symmetric $T$, that the store-convention Einstein tensor $\Ric = (\Lambda + \kappa\tau)g - 2\kappa T$ has own-ray classes $\kappa T(\ell,\ell)$ on eleven null vectors, trace $4\Lambda + 2\kappa\tau$, value $\Lambda g$ at $T = 0$, and that the linear system of the Bianchi step forces $dh = \kappa\,d\tau$ (and, at $T = 0$, $dR = 0$).
